\documentclass[10pt,a4paper]{amsart}
\usepackage[margin=19mm]{geometry}
\usepackage{amsmath,amssymb,mathtools}
\usepackage[scr=rsfs,cal=euler]{mathalfa}
\usepackage{xfrac}
\usepackage[unicode,hidelinks,bookmarks=false]{hyperref}
\newcommand{\ie}{i.e.\ }
\newcommand{\cf}{cf.\ }
\newcommand{\resp}{respectively\ }
\newcommand{\Subsection}{\subsection}
\providecommand{\nobreakdash}{\nobreak\hbox{-}}
\setlength{\emergencystretch}{2em}
\multlinegap0pt
\usepackage{bm}
\usepackage{amssymb}
\usepackage{enumerate}
\usepackage{tikz}
\usepackage[normalem]{ulem}
\usepackage{xcolor}
\newtheorem{prop}{Proposition}
\newtheorem{thm}{Theorem}
\newcommand{\C}{\mathbb C}
\newcommand{\D}{\mathbb{D}}
\newcommand{\abs}[1]{\left|#1\right|}
\newcommand{\brkt}[1]{\left(#1\right)}
\renewcommand{\d}{\partial}
\newcommand{\set}[1]{\left\{#1\right\}}
\title{Zeros of Laurent multiple orthogonal polynomials on the unit circle}
\author{Rostyslav Kozhan, Marcus Vaktn\"as}
\date{Source: arXiv:2603.21468v1 (CC BY 4.0)}
\begin{document}
\maketitle

\noindent\emph{Source and licence.} Zeros of Laurent multiple orthogonal polynomials on the unit circle. Version arXiv:2603.21468v1 (CC BY 4.0), distributed by arXiv under the Creative Commons Attribution 4.0 International licence, http://creativecommons.org/licenses/by/4.0/, as recorded in the arXiv licence metadata for this exact version and shown on the abstract page https://arxiv.org/abs/2603.21468v1. Full licence notice: \texttt{licenses/arxiv-2603.21468-CC-BY-4.0.txt}.\par\smallskip

\noindent\textit{Editorial scope.} Counted unit: the article's thm polynomials (source0.tex lines 760--762). The article's complete original proof of it is delivered with it (source0.tex lines 763--786, one \texttt{proof} environment of 24 lines). No mathematics is written, repaired or re-derived: the statement and the proof are the article's own lines, in the article's own notation, and no retained line is substituted or re-flowed.  Retained uncounted as the source-local context the proof needs: lines 228--231; lines 233--235; lines 449--452; lines 455--457; lines 459--467; lines 492--495; lines 510--512; lines 514--516; lines 523--525; lines 711--713.  Nothing was omitted from any retained span, so the package carries no omission record.  No retained line contains a citation, so the wrapper carries no bibliography and no bibliography entry.  No retained line contains a figure environment, an image-inclusion command or drawing code, so no image is delivered and the package declares no asset dependency.  Editorial formatting only, in addition: an AMS \texttt{amsart} preamble that restates the article's own declarations and macros used by the retained lines, namely the environments \texttt{prop}, \texttt{thm} on separate counters, together with the standard packages those lines need.  The retained environments print in this document as Proposition 1, Theorem 1 in order of appearance, whereas the article prints the same items with the numbering of its own full document.  The complete native source of the article as distributed in the arXiv e-print for this exact version is bundled unchanged as \texttt{sources/196950/source0.tex}.
    \begin{equation}
\label{eq:generalized hermite pade orthogonality intro}
    \int \Phi_{\bm{n},\bm{m}}(w)w^{-p}d\mu_j(w) = 0, \quad p = -m_j,\dots,n_j-1, \quad j = 1,\dots,r.
\end{equation}

\begin{equation}\label{eq:moprlIIspan}
    \Phi_{\bm{n},\bm{m}}(z) = z^{\abs{\bm{n}}} + \ldots + \alpha_{\bm{n},\bm{m}}z^{-\abs{\bm{m}}}.
    \end{equation} 

    \begin{equation}
\label{eq:generalized hermite pade orthogonality star}
    \int \Phi_{\bm{n},\bm{m}}^*(w)w^{-p}d\mu_j(w) = 0, \quad p = -m_j+1,\dots,n_j, \quad j = 1,\dots,r.
\end{equation}

\begin{equation}\label{eq:moprlIIspanStar}
    \Phi_{\bm{n},\bm{m}}^*(z) = \beta_{\bm{n},\bm{m}}z^{\abs{\bm{n}}} + \ldots + z^{-\abs{\bm{m}}}.
\end{equation}

\begin{prop}\label{prop:normality criteria hp}
    Normality of $(\bm{n},\bm{m})$ is equivalent to the following statements.
    \begin{itemize}
        \item There is a unique $\Phi_{\bm{n},\bm{m}}$ on the form \eqref{eq:moprlIIspan} satisfying \eqref{eq:generalized hermite pade orthogonality intro}.
        \item There is no $\Phi_{\bm{n},\bm{m}}(z) \in \operatorname{span}\{z^p\}_{p = -\abs{\bm{m}}}^{\abs{\bm{n}}-1}\setminus\{0\}$ satisfying \eqref{eq:generalized hermite pade orthogonality intro}.
        \item There is a unique $\Phi_{\bm{n},\bm{m}}^*$ on the form \eqref{eq:moprlIIspanStar} satisfying \eqref{eq:generalized hermite pade orthogonality star}.
        \item There is no $\Phi_{\bm{n},\bm{m}}^* \in \operatorname{span}\{z^p\}_{p = -\abs{\bm{m}}+1}^{\abs{\bm{n}}}\setminus\{0\}$ satisfying \eqref{eq:generalized hermite pade orthogonality star}.
    \end{itemize}
\end{prop}

\begin{equation}\label{eq:paraorthogonal polynomials defn}
    % X_{\bm{n}}^{(\tau)}(z) = \tau^{1/2} z^{1/2}\phi_{\bm{n}}(z) + \tau^{-1/2} z^{-1/2}\phi_{\bm{n}}^\sharp(z).
    X_{\bm{n}}^{(\tau)}(z) = z^{1/2}\phi_{\bm{n}}(z) + \tau z^{-1/2}\phi_{\bm{n}}^\sharp(z).
\end{equation}

\begin{equation}\label{eq:orthoPOPUC}
    \int X_{\bm{n}}^{(\tau)}(z)z^{-p}d\mu_j(z) = 0, \quad p = -(n_j-1)/2,\dots,(n_j-1)/2.
\end{equation}

\begin{equation}\label{eq:beta invariance}
    X_{\bm{n}}^{(\tau)}(z) = \tau\overline{X_{\bm{n}}^{(\tau)}(1/\bar{z})},
\end{equation}

\begin{prop}\label{prop:popuc normality}
    $\bm{n}$ is $\phi$-normal if and only if there is no non-trivial $\tau$-invariant $X(z) \in \operatorname{span}\set{z^p}_{p = -(\abs{\bm{n}}-1)/2}^{(\abs{\bm{n}}-1)/2}$ satisfying \eqref{eq:orthoPOPUC}. 
\end{prop}

\begin{thm}\label{thm:zeros in unit disc}
    For Angelesco and AT systems, the polynomial $z^{\abs{\bm{n}}/2}\phi_{\bm{n}}(z)$ has $\abs{\bm{n}}$ zeros in $\D$, with each zero counted as many times as its multiplicity. 
\end{thm}

\begin{thm}\label{thm:normality of hp polynomials}
    Let $\bm{\mu}$ be an Angelesco system or an AT system. Then the multi-indices $(\bm{n},\bm{n}+\bm{e}_j)$ and $(\bm{n}+\bm{e}_j,\bm{n})$ are normal for each $j = 1,\dots,r$. %Furthermore, $\Phi_{\bm{n},\bm{n}+\bm{e}_j}$
\end{thm}

\begin{proof}
    By Proposition \ref{prop:normality criteria hp} it suffices to show that the polynomial $P(z) = z^{\abs{\bm{n}}+1}\Phi_{\bm{n},\bm{n}+\bm{e}_j}(z)$ has $2\abs{\bm{n}}+1$ zeros (with each zero counted according to its multiplicity), given the orthogonality relations and degree restrictions of $\Phi_{\bm{n},\bm{n}+\bm{e}_j} \not\equiv 0$. Note that $\Phi(z) = z^{1/2}\Phi_{\bm{n},\bm{n}+\bm{e}_j}(z)$ satisfies the orthogonality relations
    \begin{equation}
        \int \Phi(z)z^{-p}d\mu_j(z) = 0, \quad p = -n_j+1/2,\dots,n_j-1/2, \quad j = 1,\dots,r. 
    \end{equation}
    
    Now consider the function
    \begin{equation}
        X_{2\bm{n}}^{(\tau)}(z) = z^{1/2}\Phi_{\bm{n},\bm{n}+\bm{e}_j}(z) + \tau z^{-1/2}\Phi_{\bm{n},\bm{n}+\bm{e}_j}^\sharp(z). 
    \end{equation}
    Observe that $X^{(\tau)}_{2\bm{n}}(z)\in\operatorname{span}\{z^p\}_{p=-|\bm{n}|-1/2}^{|\bm{n}|+1/2}$ and satisfies \eqref{eq:orthoPOPUC}--\eqref{eq:beta invariance} for index $2\bm{n}$. Note that its highest coefficient is not $1$, unlike the polynomial defined in~\eqref{eq:paraorthogonal polynomials defn}.
    
    By Proposition \ref{prop:popuc normality}, the polynomial $P^{(\tau)}(z) = z^{\abs{n}+1/2}X_{2\bm{n}}^{(\tau)}(z)$ has degree $2\abs{\bm{n}}+1$. We also have
    \begin{equation}
        P^{(\tau)}(z) = P(z) + \tau P^*(z),
    \end{equation}
    where $P^*(z) = z^{2\abs{\bm{n}}+1}\overline{P(1/\bar{z})}$. From here we can repeat the ideas from the proof of Theorem \ref{thm:zeros in unit disc}.
    
    First, it follows similarly that $P$ cannot have zeros on $\d\D$, and then we can introduce the Blaschke product $B(z) = P(z)/P^*(z)$. We have $B(e^{i\theta}) = e^{i\Psi(\theta)}$ for the continuous phase function 
    \begin{equation}
        \Psi(\theta) = 2\sum_{j = 1}^{N}\brkt{-\frac{\theta}{2} + \operatorname{Arg}(e^{i\theta}-z_j)}, \quad \theta \in [-\pi,\pi),
    \end{equation}
    where $N = \deg{P} \leq 2\abs{\bm{n}}+1$ and $z_1,\dots,z_N$ are the zeros of $P$. The argument principle then gives $\Psi(\pi-)-\Psi(-\pi) = 2\pi(n_+-n_-) = \pm2\pi(2\abs{\bm{n}}+1)$, where the sign depends on whether $\Psi$ is increasing or decreasing. Depending on the sign we then have either $2\abs{\bm{n}}+1$ zeros in $\D$ or $2\abs{\bm{n}}+1$ zeros in $\C\setminus\bar\D$. Either case proves $\deg{P} = 2\abs{\bm{n}}+1$, which completes the proof. 
    \end{proof}

\end{document}
